\setcounter{section}{1}

\Needspace{5\baselineskip}
\hypertarget{dpo}{%
\section{DPO}\label{dpo}}

DPO is commonly used for RLHF.

Its input is a dataset of triples \(\mathcal{D}=\{x,y_w,y_l\}\).

\Needspace{5\baselineskip}
Model preparation includes:

\begin{itemize}
\tightlist
\item
  \(\pi_\theta\) (policy model): the model to train, initialized with SFT weights.
\item
  \(\pi_{ref}\) (reference model): usually a frozen copy of the SFT model, excluded from parameter updates.
\end{itemize}

\Needspace{5\baselineskip}
In RLHF, the goal is to maximize expected reward while using KL divergence to prevent excessive deviation from the reference (SFT) model. The objective is:

\[
\max_{\pi}\mathbb{E}_{x\sim D,y\sim\pi(\cdot|x)}[r^*(x,y)]-\beta D_{\mathrm{KL}}(\pi(\cdot|x)\Vert\pi_{ref}(\cdot|x))
\]

\Needspace{5\baselineskip}
Here:

\begin{itemize}
\tightlist
\item
  \(x\): the prompt.
\item
  \(y\): the response.
\item
  \(r^*(x,y)\): the true, unknown reward function.
\item
  \(\pi_{ref}\): the reference model, usually the SFT model.
\end{itemize}

\Needspace{5\baselineskip}
It can be proved that this problem's optimal solution \(\pi^*\) has a closed form:

\[
\pi^*(y|x)=\frac{1}{Z(x)}\pi_{ref}(y|x)\exp\left(\frac{1}{\beta}r^*(x,y)\right)
\]

Here, \(Z(x)=\sum_y\pi_{ref}(y|x)\exp\left(\frac{1}{\beta}r^*(x,y)\right)\) is the partition function.

\Needspace{5\baselineskip}
The key step in DPO rearranges this equation to solve for the reward function \(r^*(x,y)\):

\[
r^*(x,y)=\beta\log\frac{\pi^*(y|x)}{\pi_{ref}(y|x)}+\beta\log Z(x)
\]

Thus, knowing the optimal policy \(\pi^*\) and reference policy \(\pi_{ref}\) allows the reward to be expressed through the logarithm of their ratio.

\Needspace{5\baselineskip}
Human preference data usually contains an input \(x\) and two responses \((y_w,y_l)\), with \(y_w\) the winning (better) response and \(y_l\) the losing one. The Bradley--Terry model defines the probability of selecting \(y_w\):

\[
P(y_w>y_l|x)=\sigma(r^*(x,y_w)-r^*(x,y_l))
\]

Here, \(\sigma\) is Sigmoid.

\Needspace{5\baselineskip}
Express the reward difference through policy ratios:

\[
r^*(x,y_w)-r^*(x,y_l)=\beta\log\frac{\pi^*(y_w|x)}{\pi_{ref}(y_w|x)}-\beta\log\frac{\pi^*(y_l|x)}{\pi_{ref}(y_l|x)}
\]

\Needspace{5\baselineskip}
Finally, use parameterized model \(\pi_\theta\) to approximate \(\pi^*\). Substituting into the negative log-likelihood loss gives the final DPO loss:

\[
\mathcal{L}_{DPO}(\pi_\theta;\pi_{ref})=-\mathbb{E}_{(x,y_w,y_l)\sim D}\left[\log\sigma\left(\beta\log\frac{\pi_\theta(y_w|x)}{\pi_{ref}(y_w|x)}-\beta\log\frac{\pi_\theta(y_l|x)}{\pi_{ref}(y_l|x)}\right)\right]
\]

Intuitively, increase the probability ratio for \(y_w\) (the good response) relative to the reference model, and decrease it for \(y_l\) (the bad response). \(\beta\) controls how far the model deviates from the reference.

The larger the difference between these terms, the higher the probability after Sigmoid. Taking the negative logarithm is equivalent to cross-entropy loss, pushing the predicted probability toward 1.

\FloatBarrier
\Needspace{5\baselineskip}
\hypertarget{references-80}{%
\subsection{References}\label{references-80}}

\vspace{0.25em}
\begin{itemize}
\tightlist
\item
  Rafailov, R., Sharma, A., Mitchell, E., Manning, C. D., Ermon, S., \& Finn, C. (2023). \href{https://arxiv.org/abs/2305.18290}{Direct Preference Optimization: Your Language Model is Secretly a Reward Model}. NeurIPS 2023.
\item
  Ouyang, L., Wu, J., Jiang, X., et al.~(2022). \href{https://arxiv.org/abs/2203.02155}{Training Language Models to Follow Instructions with Human Feedback}. NeurIPS 2022.
\end{itemize}

\clearpage
